\documentclass[12pt]{article}

\usepackage{bbm}
\usepackage{fancyhdr}
\usepackage{titling}
\usepackage{enumitem}
\usepackage{amsmath}
\usepackage{dsfont}
\usepackage{amsfonts}
\usepackage{amssymb}
\usepackage{amsthm}
\usepackage{geometry}
\usepackage{tikz}
\usepackage{algorithm}
\usepackage{algpseudocode}
\geometry{verbose,tmargin=1in,bmargin=1in,lmargin=1in,rmargin=1in}
\renewcommand{\vec}[1]{\mathbf{#1}}
\newcommand{\dec}{\text{dec}}
\newcommand{\divi}{\mathbin{\textnormal{div}}}
\title{Homework 1}
\author{Satvik Prasad}
\date{}
\newtheorem{lemma}{Lemma}[section]
\fancypagestyle{plain}{%
  \fancyhf{}
  \fancyhead[L]{\theauthor}
  \fancyhead[R]{\thedate}
}
\pagestyle{plain}
\makeatletter
\def\@maketitle{%
  \newpage
  \null
  \vskip 1em%
  \begin{center}%
    \let \footnote \thanks
    {\LARGE \@title \par}%
    \vskip 1em%
  \end{center}%
  \par
\vskip 1em}
\makeatother
\usepackage{cmbright}

\newcommand{\mat}[1]{\mathbf{#1}}

% ---------- Probability theory ----------
% Core operators (optional argument-free; use as \E{X}, \Var{X}, \Cov{X}{Y})
\DeclareMathOperator{\EOp}{\mathbb{E}}
\DeclareMathOperator{\VarOp}{Var}
\DeclareMathOperator{\CovOp}{Cov}
\DeclareMathOperator{\CorrOp}{Corr}

\renewcommand{\P}[1]{\mathbb{P}\left(#1\right)}        % \P{A}
\newcommand{\E}[1]{\EOp\left[#1\right]}                % \E{X}
\newcommand{\Var}[1]{\VarOp\left(#1\right)}            % \Var{X}
\newcommand{\Cov}[2]{\CovOp\left(#1,\,#2\right)}       % \Cov{X}{Y}
\newcommand{\Corr}[2]{\CorrOp\left(#1,\,#2\right)}     % \Corr{X}{Y}
\newcommand{\SD}[1]{\operatorname{SD}\left(#1\right)}  % \SD{X}

% Subscripted versions, e.g. \Esub{\theta}{X}, \Psub{\theta}{A}
\newcommand{\Esub}[2]{\EOp_{#1}\left[#2\right]}
\newcommand{\Psub}[2]{\mathbb{P}_{#1}\left(#2\right)}

% Conditioning, indicators, independence
\newcommand{\given}{\,\middle|\,}                      % \E{X \given Y}
\newcommand{\ind}[1]{\mathbbm{1}\left\{#1\right\}}     % \ind{X > 0}
\newcommand{\indep}{\perp\!\!\!\perp}                  % X \indep Y
\newcommand{\iid}{\overset{\text{iid}}{\sim}}          % X_i \iid \Normal(0,1)

% Convergence
\newcommand{\convp}{\xrightarrow{\;p\;}}               % in probability
\newcommand{\convd}{\xrightarrow{\;d\;}}               % in distribution
\newcommand{\convas}{\xrightarrow{\text{a.s.}}}        % almost surely
\newcommand{\convLp}[1]{\xrightarrow{L^{#1}}}          % in L^p

% Probability spaces
\newcommand{\F}{\mathcal{F}}                           % sigma-algebra
\newcommand{\B}{\mathcal{B}}                           % Borel sigma-algebra
\newcommand{\pspace}{(\Omega, \F, \mathbb{P})}

% Common distributions
\newcommand{\Normal}{\mathcal{N}}
\newcommand{\Bern}{\operatorname{Bernoulli}}
\newcommand{\Bin}{\operatorname{Binomial}}
\newcommand{\Pois}{\operatorname{Poisson}}
\newcommand{\Unif}{\operatorname{Uniform}}
\newcommand{\Exp}{\operatorname{Exp}}
\newcommand{\Geom}{\operatorname{Geometric}}
\newcommand{\Gam}{\operatorname{Gamma}}
\newcommand{\Beta}{\operatorname{Beta}}

% Number sets
\newcommand{\R}{\mathbb{R}}
\newcommand{\N}{\mathbb{N}}
\newcommand{\Z}{\mathbb{Z}}


\begin{document}
\begin{enumerate}
    \item We have some weighted digraph $G$ such that every shortest path from $s$ has at most $k$ edges. We can continuously relax edges, until we reach an iteration where no edge is relaxed.

    \begin{algorithm}
        \begin{algorithmic}
            \Procedure{ShortestKPath}{$G, s$}
                \State $\texttt{dist} := \texttt{map}(\infty)$ \Comment{Map where every vertex has default entry $\infty$}
                \State $\texttt{dist}[s] \gets 0$
                \State $\texttt{changed} \gets \texttt{true}$

                \vspace{3mm}

                \While{\texttt{changed}} 
                    \State $\texttt{changed} \gets \texttt{false}$

                    \vspace{3mm}
                        
                    \For{$(u, v) \in \mathbb E(G)$}
                        \If{$\texttt{dist[u]} + c(u, v) < \texttt{dist}[v]$}
                            \State $\texttt{dist}[v] \gets \texttt{dist}[u] + c(u, v)$
                            \State $\texttt{changed} \gets \texttt{true}$
                        \EndIf
                    \EndFor
                \EndWhile

                \vspace{3mm}

                \State \Return \texttt{dist}
            \EndProcedure
        \end{algorithmic}
    \end{algorithm}

    Now, we prove that such an algorithm is $\mathcal O(k |E|).$ Every iteration, we relax $\mathcal O(|E|)$ edges, so it suffices to prove that we have at most $k$ iterations. We prove this via induction.

    Let us assume that at iteration $j \in \mathbb Z_{\geq 0},$ $\texttt {dist}[v] = \text{dist}(s, v)$ for all vertices $v$ whose shortest path has $\leq j$ edges. 
    
    Then, consider any vertex $w$ whose shortest $s-w$ path has $\leq j + 1$ edges. Then, let $w'$ be the predecessor of $w$ in this path. If the $s-w'$ path is not the shortest possible path, then the $s-w$ path is not the shortest (we can form one shorter), so $s-w'$ path must be the shortest path, and take $j$ edges. Hence, $\texttt{dist}[w'] = \text{dist}(s, w'),$ meaning our update rule will eventually relax $\texttt{dist}[w]$ to $\text{dist}(s, w') + c(w', w) = \text{dist}(s, w),$ since the $s-w$ path is minimal.

    On iteration $0,$ all vertices with shortest path with $0$ edges from $s$ (i.e., only $s$) have their distance correctly set (to $0$). Hence, at iteration $j,$ the $\texttt{dist}$ map is correct up to vertices with shortest path at most $j$ length, for all $j \geq 0.$

    Since we terminate the loop when our relaxation never changes edge weights, we terminate at iteration $i,$ where $i$ is the smallest edge length such that there are no shortest paths from $s$ to any vertex $v$ of length $i.$ Thus, $i \leq k + 1,$ meaning we have $\mathcal O(k)$ iterations.

    Thus, the algorithm is correct and runs in $\mathcal O(k |E|)$ time.

    \newpage

    \item We can use dynamic programming but also store this new auxiliary table, given some digraph $G.$

    \begin{algorithm}
        \caption{Modified Floyd Warshall}
        \begin{algorithmic}
            \Procedure{ModifiedFloydWarshall}{$G$}
                \State $\texttt{dist}[i][j] = 0$ when $i = j$ and $\infty$ otherwise
                \State $\texttt{le}[i][j] = \texttt{nil}$ for all $i, j$
                \State $n = |V(G)|$

                \vspace{3mm}

                \For{$i \in 2 \ldots n$}
                    \For{$s \in 1 \ldots i - 1$}
                    \Comment{Find shortest path to and from $i$}
                        \For{$t \in 1 \ldots i - 1$}
                            \If{$\texttt{dist[s][t]} + c(t, i) < \texttt{dist}[s][i]$}
                                \State $\texttt{dist}[s][i] \gets \texttt{dist}[s][t] + c(t, i)$
                                \State $\texttt{le}[s][i] \gets (t, i)$
                            \EndIf

                            \vspace{3mm}
                            
                            \If{$\texttt{dist[t][s]} + c(i, t) < \texttt{dist}[i][s]$}
                                \State $\texttt{dist}[i][s] \gets \texttt{dist}[t][s] + c(i, t)$
                                \State $\texttt{le}[i][s] \gets \texttt{le}[t][s]$
                            \EndIf
                        \EndFor

                        \vspace{3mm}
                    \EndFor
                    
                    \vspace{3mm}

                    \For{$s \in 1 \ldots i - 1$}
                    \Comment{Relax using intermediate as $i$}
                        \For{$t \in 1 \ldots i - 1$}
                            \If{$\texttt{dist}[s][i] + \texttt{dist}[i][t] < \texttt{dist}[s][t]$}
                                \State $\texttt{dist}[s][t] \gets \texttt{dist}[s][i] + \texttt{dist}[i][t]$
                                \State $\texttt{le}[s][t] = \texttt{le}[i][t]$
                            \EndIf
                        \EndFor
                    \EndFor
                \EndFor

                \vspace{3mm}

                \State \Return \texttt{dist}, \texttt{le}
            \EndProcedure
        \end{algorithmic}
    \end{algorithm}

    \newpage

    \item We biject this scenario to a graph problem. Assume we are given $n, m$ and $I[1\ldots m],$ $J[1 \ldots m]$ and $C[1 \ldots m].$ Construct digraph $G$ with vertices $1, 2, \ldots, n,$ and weighted edge $(i, j)$ if and only if there is some index $k$ such that $I[k] = i, J[k] = j,$ with weight $C[k].$ Equivalently, 
    $$
    E(G) = \{(I[k], J[k], C[k]) : 1 \leq k \leq m\}.
    $$

    We prove that the set of equations has a solution if and only if the graph has no negative weight cycle.

    \textbf{Forward}: Assume the graph has a negative edge cycle. Let a traversal of this cycle be $i_1, i_2, \ldots, i_k$ for $1 \leq k \leq n,$ with edge weights $c_{i_1}, \ldots, c_{i_k}.$ Assuming the set of equations has a solution, then, this cycle corresponds to the equations
    \begin{align*}
    x_{i_1} &\leq  x_{i_2} + c_{i_1} \\
    x_{i_2} &\leq  x_{i_3} + c_{i_2} \\
            &\vdots \\
    x_{i_k} &\leq  x_{i_1} + c_{i_k},
    \end{align*}
    and after summing all equations we get
    $$
        x_{i_1} + \ldots + x_{i_k} \leq x_{i_1} + \ldots + x_{i_k } + \sum_{j = 1}^{k} c_{i_j} \Rightarrow \sum_{j = 1}^{k}c_{i_j} \geq 0,
    $$
    which is the weight of our cycle. However, the cycle has a negative weight, meaning this is a contradiction. Thus, our graph having a negative weight cycle implies our set of equations has no solution. By contraposition, if our set of equations has a solution, there is no negative weight cycle.

    \textbf{Reverse:} Assume that the graph has no negative weight cycle. Then, root the graph at vertex $1,$ and set $x_1 = 0,$ and let $x_j = \text{dist}(j, 1),$ where $\text{dist}(i , j)$ is the shortest path distance from $i$ to $j.$ Since there are no negative weight cycles, such a shortest path distance is guaranteed to exist.

    Then, consider any path from $i \to j,$ and their corresponding inequality $$x_i \leq x_j + c(i \to j).$$ Assume this inequality is violated for an arbitrary pair $i, j.$

    Since $x_i = \text{dist}(i, 1)$ and $x_j = \text{dist}(j, 1),$ we get
    $$
    \text{dist}(i, 1) > \text{dist}(j, 1) + c(i \to j)
    $$
    However, note that $\text{dist}(j, 1) + c(i \to j)$ is the weight of a path from $i \to 1$ through $j$ that is more minimal than $\text{dist}(i, 1).$ This is a contradiction, meaning the inequality holds for all possible pairs $i, j,$ and thus we have a valid solution.

    Therefore, there is a solution if and only if the associated digraph has no negative weight cycle. Thus, a valid solution can run Bellman-Ford on digraph $G,$ and if we abort due to relaxations carrying on past the $n$-th iteration, we return $\texttt{NO},$ and return $\texttt{YES}$ otherwise.
\end{enumerate}
\end{document}